\subsection{Regional canonical reconstruction from two field layers}
\label{sec:source-scalar-regional-time-slice}

The scalar action permits a concrete regional algebra construction. Its
access requirement is stronger than a stored classical configuration:
one must admit the full real-parameter Weyl families of the specified
linear field observables. No quantum outcomes or physical region assignment
follow from a classical trace.

Let $M=\operatorname{diag}(m_i)>0$, let $K=K^T$ be the complete stiffness
matrix, and put $A=M^{-1}K$. Take $\hbar>0$. In the finite Schr\"odinger representation,
$[q_i,p_j]=i\hbar\delta_{ij}$, and the position component of one canonical
kick--drift--kick step of duration $\tau\ne0$ is
\[
 q^+=q+\tau M^{-1}p-\frac{\tau^2}{2}Aq.
\]
This is the full coupled action, including the entries connecting a region
to its complement. For a site set $R$, write $E=R^c$, $r=|R|$,
$B=A_{R,E}$, and let $Cq_E$ be the admitted exterior linear fields.
For a real space $S$ of linear canonical observables, write
$\mathcal W(S)=\{\exp(i\ell/\hbar):\ell\in S\}^{\prime\prime}$.

\begin{theorem}[Sharp regional collar criterion]
\label{thm:source-scalar-regional-collar}
Let $\mathcal A_R=\mathcal W(\operatorname{span}_{\mathbb R}\{q_i,p_i:i\in R\})$
and
\[
 \mathcal D_R(C)=\mathcal W\bigl(
    \operatorname{span}_{\mathbb R}\{q_R,q_R^+,Cq_E\}\bigr).
\]
Then
\[
 \mathcal A_R\subseteq\mathcal D_R(C)
 \quad\Longleftrightarrow\quad
 \ker C\subseteq\ker B
 \quad\Longleftrightarrow\quad
 \operatorname{row}B\subseteq\operatorname{row}C.
\]
Consequently the minimum number of independent exterior linear fields is
$\operatorname{rank}B$. Choosing a row basis $C$ of $B$ attains this minimum.
With no exterior access the original regional algebra is recovered exactly
if and only if $B=0$.
\end{theorem}

\begin{proof}
Using the local $q_R$ directions and the invertibility of $M_R$,
the real linear space defining $\mathcal D_R(C)$ is equivalently generated by
\[
 q_R,\qquad p_R-\frac{\tau}{2}M_RBq_E,\qquad Cq_E.
\]
Every $p_R$ direction is therefore present exactly when every row of $B$
is in the row space of $C$. The finite-dimensional kernel criterion and
the sharp rank bound follow by annihilators. Sufficiency also gives an
explicit reconstruction: if $B=DC$, then
\[
 p_R=M_R\left[
   \frac{q_R^+-q_R}{\tau}
   +\frac{\tau}{2}\bigl(A_{R,R}q_R+DCq_E\bigr)\right].
\]
For necessity, if $Ce=0$ but $Be\ne0$, the perturbations
$\delta q_R=0$, $\delta q_E=e$ and
$\delta p_R=(\tau/2)M_RBe$ leave all admitted linear data unchanged
but change a local momentum.

The passage from linear spans to von Neumann algebras does not assume
that observable products are linear. If a linear canonical direction $v$
lies outside a real subspace $S$, nondegeneracy of the ambient symplectic
form gives $z\in S^\sigma$ with $\sigma(z,v)\ne0$. Every Weyl operator
in direction $z$ commutes with $\mathcal W(S)$, whereas a suitable real
multiple does not commute with the Weyl family in direction $v$.
That family cannot lie in $\mathcal W(S)$. Conversely the Weyl relations
produce every direction in the real span. This proves the algebraic
equivalence, including degenerate $S$.
\end{proof}

The two field layers themselves are noncommutative:
\[
 [q_i,q_j^+]=i\hbar\frac{\tau}{m_i}\delta_{ij},
 \qquad [q_i^+,q_j^+]=0.
\]
Their $2r$-dimensional regional space is nondegenerate, even when it fails
to contain the original local momenta. In fact
\[
 \frac{m_i}{\tau}q_i^+=p_i+(Gq)_i,
 \qquad G=M/\tau-\tau K/2=G^T.
\]
Multiplication by
$U=\exp[-iq^TGq/(2\hbar)]$ implements
$UpU^*=p+Gq$. Hence
$\mathcal D_R(0)=U\mathcal A_RU^*$: the two-layer family is an isotonic
finite noncommutative net, and its disjoint site regions commute. The same
unitary works for every region. This sheared net is not identified with
the original site net when exterior couplings are present. For a minimal
collar, the enlarged real space instead has dimension
$2r+\operatorname{rank}B$ and a central exterior subspace of dimension
$\operatorname{rank}B$.

For the same 64-site golden action, exact arithmetic gives the following
minimal access counts. The coordinate column counts every exterior site
with a nonzero incident coupling; a linear combination can require reads
from several such sites.
\[
\begin{array}{lrrr}
\text{region}&r&\text{exterior coordinate fields}&\operatorname{rank}B\\
\hline
\text{interior singleton}&1&6&1\\
\text{central eight-site cube}&8&24&8\\
\text{half-volume}&32&16&16\\
\text{interior four-site line}&4&16&4\\
\text{checkerboard}&32&32&32\\
\text{whole volume}&64&0&0
\end{array}
\]
The complete certificate also retains the empty region and corner
singleton. It verifies both exact rank factorizations, an identity minor,
the full cross-time canonical brackets, and 62 explicit perturbations
showing that deletion of any selected collar row loses a local momentum.
These are mathematical access counts, not implemented routing costs.

The kernel criterion and explicit velocity reconstruction have Lean
proofs; the Weyl and von Neumann interpretation above is analytic.
Native availability of the continuous Weyl families, preparation and
outcome maps, physical clock and region calibration, and compatibility
with a continuum regional net remain additional requirements.
